% Input after (ST.1)--(ST.3) and the cone-probability notation.
\subsection{三次端点的全局相对 \texorpdfstring{$L^1$}{L1} 稳定性}
\label{subsec:global-l1-cubic}

这是给定变差 $f$-散度框架及勘误
\cite{Gilardoni2006MinimumF,Gilardoni2010Corrigendum}下的显式特例；以下自足优化二点边界。

\begin{proposition}[三次矩的全局常数]\label{prop:global-l1-cubic}
令 $(\Omega,\mathcal A,\mathbb P)$ 为任意概率空间，$f\ge0$ 且
$\mathbb Ef=1$，并置
\[
 d=\mathbb E|f-1|,\qquad D_3=\mathbb Ef^3-1.
\]
则
\[
 D_3\ge C_3d^2,\qquad C_3=\frac{11+5\sqrt5}{8}.\tag{GL.1}
\]
除 $f=1$ 外，等号成立当且仅当存在概率为
\(q_*=(3-\sqrt5)/2\) 的集合 $A$，使得 $f=0$ 于 $A$，而在补集上
\(f=(1-q_*)^{-1}\)。

回到 \textup{(ST.1)} 并令 $p=3$（等价于 $\tau=2n/3$），有
\[
 d(h)=\frac{\int_K|h-h_{\rm rad}|}{\int_Kh}
 \le\sqrt{\frac{S^3-1}{C_3}}.\tag{GL.2}
\]
当 $n\ge2$ 时，\textup{(GL.2)} 在满足 \textup{(H.1)} 的同次权类中尖锐。
当 $n=1$ 时，最佳常数为 $3$，且 $D_3=3d(h)^2$。
因此 \textup{(GL.2)} 将 \textup{(ST.3)} 在 $p=3$ 时的分母 $2$ 提高为 $C_3$；
但小亏损极限仍以 \textup{(ST.4)} 的系数 $1$ 更精确地描述。
\end{proposition}
\begin{proof}
只需证明 \textup{(GL.1)}。若 $\mathbb Ef^3=\infty$，结论显然；以下设该矩有限。
若 $d=0$ 则 $f=1$ a.e.。以下设 $d>0$，令
\[
 A=\{f<1\},\qquad q=\mathbb P(A),\qquad
 \delta=\mathbb E(1-f)_+=d/2.
\]
于是 $0<q<1$、$0<\delta\le q$。分别在 $A$ 与 $A^c$ 上对凸函数 $x^3$ 用
Jensen 不等式；相应条件均值为 $1-\delta/q$ 与 $1+\delta/(1-q)$。展开得
\[
 \frac{D_3}{d^2}\ge R(q,\delta):=
 \frac{3}{4q(1-q)}+
 \frac{\delta(2q-1)}{4q^2(1-q)^2}.\tag{GL.3}
\]
若 $q\ge1/2$，右端不小于 $3$。若 $0<q<1/2$，右端随 $\delta$ 递减，故
\[
 R(q,\delta)\ge R(q,q)=G(q):=\frac{2-q}{4q(1-q)^2}.
\]
此处
\[
 G'(q)=\frac{(q-1)(q^2-3q+1)}{2q^2(1-q)^4}.
\]
因此 $G$ 在 $(0,1/2)$ 上先减后增，其唯一极小点是
\(q_*=(3-\sqrt5)/2\)，且直接代入给出 $G(q_*)=C_3<3$。这证明
\textup{(GL.1)}。非平凡等号迫使上述每一步取等：
$q=q_*$、$\delta=q_*$，以及两个条件 Jensen 不等式取等，恰给出命题所列二值函数。

由 \textup{(ST.1)}，当 $p=3$ 时 $S^3=\mathbb Ef^3$ 且
$d(h)=\mathbb E|f-1|$，故 \textup{(GL.2)} 已证。为说明 $n\ge2$ 的尖锐性，
令 $\pi(x)=x/\rho(x)$（$x\in K\setminus\{0\}$）。锥测度是归一 Lebesgue 测度经
$\pi$ 的推前；每个 $\pi^{-1}\{\theta\}$ 是一条线段，Lebesgue 测度为零。由于
$\partial K$ 是紧度量空间，$\sigma_K$ 因而无原子，故可取
$\sigma_K(A)=q_*$。在 $A$ 上令 $f=a$，其中 $0<a<1$，在补集上令
\(f=(1-q_*a)/(1-q_*)\)，并令 $a\downarrow0$。每个剖面对应正有界的
\textup{(H.1)} 权，且 \textup{(GL.3)} 中的比值趋于 $C_3$。

最后，$n=1$ 时锥测度为两个端点上的等质量原子。均值一迫使两个取值为
$1+\varepsilon,1-\varepsilon$（交换端点不影响计算），从而
$D_3=3\varepsilon^2$、$d=|\varepsilon|$，给出常数 $3$。
\end{proof}

\begin{remark}
该命题优化的是 $p=3$ 的首项稳定界；它不扩大任何二阶展开的参数窗口。
严格正性排除非平凡等号，却不影响 $n\ge2$ 的尖锐逼近；上一节的尖峰反例同样排除 $L^\infty$ 控制。
\end{remark}
